\section{来源审计与课程地图：同一个 attention，两种计算轴}

本讲由两个相互咬合的部分组成。Aidan Gomez 先用约十五分钟回顾 Transformer 的起源与训练直觉；随后 Jannik Kossen 与 Neil Band 介绍 Non-Parametric Transformer（NPT）。若只按演讲顺序做摘要，二者像是“基础复习加一篇论文”；真正的教学主线却更统一：Transformer 首先把 attention 放到\textbf{同一序列的 tokens 之间}，NPT 再把 attention 的作用轴提升到\textbf{同一数据集的 datapoints 之间}。

本次重写以 Stanford Online 官方上传 `zejXBg-2Vpk` 为时间基准。视频提供 1,476 条可解析的官方手动英文字幕，但没有公开独立 slide PDF，因此从 1080p 录像中恢复并人工筛选 28 张教学页。渐进动画只保留最终完整状态或真正有独立教学作用的中间状态；Q\&A 中反复跳回的旧页、空白转场与片尾被省略，口头解释则进入 teacher-voice ledger 与正文。

\lecturefigure{slide-01-transformers-title.jpg}{第一部分标题：Transformer 的简史与直觉。}{Stanford Online 官方视频 00:00:46--00:01:07。}

这张标题页提醒我们，Gomez 的任务不是重新讲完一门 Transformer 课，而是抽出三个决定后续故事的机制：self-attention、multi-head attention，以及可以并行训练的 autoregressive decoder。理解这三个机制后，读者才容易看出 NPT 的关键创新并不是“又发明一种 attention”，而是重新定义 attention 的对象。

\lecturefigure{slide-02-transformer-overview.jpg}{Transformer 概览：multi-head attention、self-attention 与 fast autoregressive decoding。}{Stanford Online 官方视频 00:01:08--00:01:41。}

\begin{importantbox}{整讲的一个压缩问题}
当模型预测一个样本时，它只能依赖“该样本的特征 + 已经压入参数的训练经验”，还是可以在推理时直接读取其他训练样本？标准 supervised neural network 通常采用前者；NPT 尝试学习后者，而且让“读哪些样本、如何组合它们”由 attention 端到端决定。
\end{importantbox}

\teachervoice{Gomez 把 Transformer 的贡献概括为三个组合件，而不是某个孤立公式。这个视角很重要：后来成为标准配置的 LayerNorm、learning-rate warmup、初始化和 residual 细节，当时都还是会显著影响成败的实验变量。}

\subsection{阅读路线}

接下来先建立 token-level attention 与 autoregressive training 的最小数学框架，再转向 NPT 的 dataset-level input contract。实验部分不会只复述“排名很好”，而会依次回答四个更强的问题：NPT 是否竞争力足够、是否能表示跨样本算法、是否真的使用其他样本、以及这种依赖具有怎样的结构。

\begin{knowledgebox}{两条轴不要混淆}
\begin{tabularx}{\textwidth}{L{0.23\textwidth}>{\raggedright\arraybackslash}X >{\raggedright\arraybackslash}X}
\toprule
机制 & 被 attention 比较的对象 & 主要问题 \\
\midrule
序列 self-attention & 一个样本内部的 tokens / attributes & 句内或样本内有哪些元素相关？ \\
NPT 的 ABD & 同一批次中的 datapoints / rows & 当前样本应从哪些其他样本读取信息？ \\
NPT 的 ABA & 一个 datapoint 内部的 attributes & 一行内部的字段应怎样变换与交互？ \\
\bottomrule
\end{tabularx}
\end{knowledgebox}

\subsection{本章小结}

本讲不是两个无关主题的拼接，而是 attention 作用域的逐级扩展。第一部分解释并行训练和序列关系；第二部分把数据集本身变成模型输入，使预测可以显式依赖其他 datapoints。

